\documentclass[10pt,a4paper]{amsart}
\usepackage[margin=1in]{geometry}
\usepackage{barnette-figs}
\usepackage{graphicx,amssymb}
\setlength{\parindent}{0pt}
\setlength{\parskip}{7pt}
\pagestyle{plain}
\begin{document}
\section*{07. The disk has one extra white face}
\begin{center}\begin{minipage}{.86\textwidth}\centering
\begin{tikzpicture}[bf,x=1.05cm,y=1.05cm]
\path[use as bounding box] (-5.80,-5.80) rectangle (5.80,5.80);
\fill[bfblack,even odd rule] (0,0) circle[radius=5.55] (0,0) circle[radius=5.0];
\fill[bfblack] (-4.330,-2.500) -- (-2.133,-0.405) -- (-2.200,3.528) -- (-4.330,-2.500) -- cycle;
\fill[bfblack] (-4.330,-2.500) -- (-0.344,-4.394) -- (-1.614,-2.366) -- (-4.330,-2.500) -- cycle;
\fill[bfblack] (-2.133,-0.405) -- (0.764,-0.167) -- (-0.963,1.243) -- (-2.133,-0.405) -- cycle;
\fill[bfblack] (-2.133,-0.405) -- (-1.614,-2.366) -- (0.900,-2.302) -- (-2.133,-0.405) -- cycle;
\fill[bfblack] (-0.344,-4.394) -- (4.330,-2.500) -- (0.900,-2.302) -- (-0.344,-4.394) -- cycle;
\fill[bfblack] (4.330,-2.500) -- (2.799,2.992) -- (2.506,-0.651) -- (4.330,-2.500) -- cycle;
\fill[bfblack] (0.000,5.000) -- (-2.200,3.528) -- (2.799,2.992) -- (0.000,5.000) -- cycle;
\fill[bfblack] (-2.200,3.528) -- (-0.963,1.243) -- (0.610,2.065) -- (-2.200,3.528) -- cycle;
\fill[bfblack] (0.764,-0.167) -- (2.799,2.992) -- (0.610,2.065) -- (0.764,-0.167) -- cycle;
\fill[bfblack] (0.764,-0.167) -- (0.900,-2.302) -- (2.506,-0.651) -- (0.764,-0.167) -- cycle;
\fill[bfR!12!bfblack] (-2.133,-0.405) -- (0.764,-0.167) -- (-0.963,1.243) -- (-2.133,-0.405) -- cycle;
\fill[bfR!12!bfblack] (-2.133,-0.405) -- (-1.614,-2.366) -- (0.900,-2.302) -- (-2.133,-0.405) -- cycle;
\fill[bfR!12!bfblack] (-2.200,3.528) -- (-0.963,1.243) -- (0.610,2.065) -- (-2.200,3.528) -- cycle;
\fill[bfR!12!bfblack] (0.764,-0.167) -- (2.799,2.992) -- (0.610,2.065) -- (0.764,-0.167) -- cycle;
\fill[bfR!12!bfblack] (0.764,-0.167) -- (0.900,-2.302) -- (2.506,-0.651) -- (0.764,-0.167) -- cycle;
\fill[bfR!7] (-2.133,-0.405) -- (-0.963,1.243) -- (-2.200,3.528) -- (-2.133,-0.405) -- cycle;
\fill[bfR!7] (-2.133,-0.405) -- (0.900,-2.302) -- (0.764,-0.167) -- (-2.133,-0.405) -- cycle;
\fill[bfR!7] (-1.614,-2.366) -- (-0.344,-4.394) -- (0.900,-2.302) -- (-1.614,-2.366) -- cycle;
\fill[bfR!7] (-2.200,3.528) -- (0.610,2.065) -- (2.799,2.992) -- (-2.200,3.528) -- cycle;
\fill[bfR!7] (-0.963,1.243) -- (0.764,-0.167) -- (0.610,2.065) -- (-0.963,1.243) -- cycle;
\fill[bfR!7] (0.764,-0.167) -- (2.506,-0.651) -- (2.799,2.992) -- (0.764,-0.167) -- cycle;
\draw[bf edge] (0,0) circle[radius=5.0];
\draw[bf edge] (-4.330,-2.500) -- (-2.133,-0.405);
\draw[bf edge] (-4.330,-2.500) -- (-1.614,-2.366);
\draw[bf edge] (-4.330,-2.500) -- (-0.344,-4.394);
\draw[bf edge] (-4.330,-2.500) -- (-2.200,3.528);
\draw[bf edge] (-2.133,-0.405) -- (-1.614,-2.366);
\draw[bf edge] (-2.133,-0.405) -- (-2.200,3.528);
\draw[bf edge] (-2.133,-0.405) -- (-0.963,1.243);
\draw[bf edge] (-2.133,-0.405) -- (0.764,-0.167);
\draw[bf edge] (-2.133,-0.405) -- (0.900,-2.302);
\draw[bf edge] (-1.614,-2.366) -- (-0.344,-4.394);
\draw[bf edge] (-1.614,-2.366) -- (0.900,-2.302);
\draw[bf edge] (-0.344,-4.394) -- (4.330,-2.500);
\draw[bf edge] (-0.344,-4.394) -- (0.900,-2.302);
\draw[bf edge] (4.330,-2.500) -- (0.900,-2.302);
\draw[bf edge] (4.330,-2.500) -- (2.506,-0.651);
\draw[bf edge] (4.330,-2.500) -- (2.799,2.992);
\draw[bf edge] (0.000,5.000) -- (-2.200,3.528);
\draw[bf edge] (0.000,5.000) -- (2.799,2.992);
\draw[bf edge] (-2.200,3.528) -- (-0.963,1.243);
\draw[bf edge] (-2.200,3.528) -- (2.799,2.992);
\draw[bf edge] (-2.200,3.528) -- (0.610,2.065);
\draw[bf edge] (-0.963,1.243) -- (0.764,-0.167);
\draw[bf edge] (-0.963,1.243) -- (0.610,2.065);
\draw[bf edge] (0.764,-0.167) -- (0.900,-2.302);
\draw[bf edge] (0.764,-0.167) -- (2.506,-0.651);
\draw[bf edge] (0.764,-0.167) -- (2.799,2.992);
\draw[bf edge] (0.764,-0.167) -- (0.610,2.065);
\draw[bf edge] (0.900,-2.302) -- (2.506,-0.651);
\draw[bf edge] (2.506,-0.651) -- (2.799,2.992);
\draw[bf edge] (2.799,2.992) -- (0.610,2.065);
\draw[bf D] (-2.133,-0.405) -- (-1.614,-2.366);
\draw[bf D] (-1.614,-2.366) -- (-0.344,-4.394);
\draw[bf D] (-0.344,-4.394) -- (0.900,-2.302);
\draw[bf D] (0.900,-2.302) -- (2.506,-0.651);
\draw[bf D] (2.506,-0.651) -- (2.799,2.992);
\draw[bf D] (2.799,2.992) -- (-2.200,3.528);
\draw[bf D] (-2.200,3.528) -- (-2.133,-0.405);
\node[] at (-0.777,0.224) {$B$};
\node[] at (-0.949,-1.691) {$B$};
\node[] at (-0.851,2.279) {$B$};
\node[] at (1.391,1.630) {$B$};
\node[] at (1.390,-1.040) {$B$};
\node[] at (-1.611,1.282) {$W$};
\node[] at (0.091,-0.946) {$W$};
\node[] at (-0.360,-3.040) {$W$};
\node[] at (0.634,2.672) {$W$};
\node[] at (0.090,1.136) {$W$};
\node[] at (1.873,0.255) {$W$};
\draw[bf root] (-4.330,-2.500) circle[radius=2.3pt];
\fill[bf vertex] (-2.133,-0.405) circle[radius=1.4pt];
\fill[bf vertex] (-1.614,-2.366) circle[radius=1.4pt];
\fill[bf vertex] (-0.344,-4.394) circle[radius=1.4pt];
\draw[bf root] (4.330,-2.500) circle[radius=2.3pt];
\draw[bf root] (0.000,5.000) circle[radius=2.3pt];
\fill[bf vertex] (-2.200,3.528) circle[radius=1.4pt];
\fill[bf vertex] (-0.963,1.243) circle[radius=1.4pt];
\fill[bf vertex] (0.764,-0.167) circle[radius=1.4pt];
\fill[bf vertex] (0.900,-2.302) circle[radius=1.4pt];
\fill[bf vertex] (2.506,-0.651) circle[radius=1.4pt];
\fill[bf vertex] (2.799,2.992) circle[radius=1.4pt];
\fill[bf vertex] (0.610,2.065) circle[radius=1.4pt];
\end{tikzpicture}
\end{minipage}\end{center}
The same disk can be read without orientation. Count the labels in the shaded region: six white faces and five black faces. The useful form of the disk identity is W equals B plus one. It says both that a cycle encloses a white face and that the surplus is rigid. When an edge of D slides onto the two spokes through its black centre, the region between the routes excludes the unchosen corner. After all slides, the enclosed spoke regions correspond to the enclosed white faces. A unit circulation around each such region therefore makes the difference between the two orientations equal six for this cycle. The other D-cycle encloses one white face. Their positive areas are the geometric reason that minimum-cycle contributions to the polynomial have a common nonzero leading coefficient.

\textbf{Status.} PROVED-HERE in both combinatorics and topology referees; independently AI-refereed, human review pending. COMPUTED white areas 6 and 1 on this pair.

\textbf{Connects to.} 4 lowest coefficient; 5 heights; 6 disk census; 8 turning.

{\footnotesize\raggedright Sources: research/\allowbreak{}2026-11-14\_astra\_referee-combinatorics/\allowbreak{}REWRITTEN-PROP-5-1.md;\allowbreak{} research/\allowbreak{}2026-11-14\_fable\_referee-topology/\allowbreak{}REFEREE.md, part b.\par}

\end{document}